\documentclass[10pt,a4paper]{amsart}
\usepackage[margin=19mm]{geometry}
\usepackage{amsmath,amssymb,mathtools}
\usepackage[scr=rsfs,cal=euler]{mathalfa}
\usepackage{xfrac}
\usepackage[unicode,hidelinks,bookmarks=false]{hyperref}
\newcommand{\ie}{i.e.\ }
\newcommand{\cf}{cf.\ }
\newcommand{\resp}{respectively\ }
\newcommand{\Subsection}{\subsection}
\providecommand{\nobreakdash}{\nobreak\hbox{-}}
\setlength{\emergencystretch}{2em}
\multlinegap0pt
\usepackage{bm}
\usepackage{bbm}
\usepackage{dsfont}
\usepackage{xspace}
\usepackage{cancel}
\usepackage{mathtools}
\usepackage{amsthm}
\newtheorem{theorem}{Theorem}
\newtheorem{fact}[theorem]{Fact}
\newtheorem{proposition}[theorem]{Proposition}
\newcommand{\kk}{\operatorname{k}}
\title[Descent along flat composed with radicalization]{Descent along flat composed with radicalization}
\author{Mohsen Asgharzadeh}
\date{Source: arXiv:2609.24577v2 (CC BY 4.0)}
\begin{document}
\maketitle

\noindent\emph{Source and licence.} Descent along flat composed with radicalization. Version arXiv:2609.24577v2 (CC BY 4.0), distributed under the Creative Commons Attribution 4.0 International licence, as recorded in the arXiv licence metadata for this exact version and shown on the abstract page https://arxiv.org/abs/2609.24577v2. Full rights notice: licenses/arxiv-2609.24577-CC-BY-4.0.txt.\par\smallskip

\noindent\textit{Editorial scope.} Counted unit: the Proposition whose source label is \texttt{None} in the article (source0.tex lines 1749--1790, source label \texttt{None}), delivered together with the article's own complete original proof of it (source0.tex lines 1792--2046, one \texttt{proof} environment of 255 lines). No mathematics is written, repaired or re-derived: statement and proof are the article's own text line for line. Required context retained uncounted: rel (lines 1043--1051). Every definition, display and label of those lines is retained unchanged. Omitted from this package and not redistributed: 1--1042, 1052--1748, 1791--1791, 2047--2095. Nothing inside a retained excerpt is omitted or altered: every retained body is delivered exactly as the article's lines are written, so no omission receipt of any kind accompanies this package. Citations: the retained excerpts carry no citation command at all, so no bibliography entry of the article is transcribed and no reference is redistributed. Reference adaptation: none was needed and none was made; every reference inside the delivered unit resolves inside it, so no unresolved reference or citation is produced. Printed slips kept literal: no printed word, symbol, hypothesis or number of the counted statement or of its proof was altered. Layout and class: the article's own class is \texttt{amsart} with packages including amscd, amssymb, amsopn, amsmath, amsthm, graphics, amsfonts, enumerate, verbatim, calc, graphicx, centernot, mathpazo; the wrapper is the lane's pinned \texttt{amsart} editorial wrapper at 10pt on A4, which restates the theorem-like environments the retained text needs and adds the article's own macro definitions that the retained text needs (\texttt{\textbackslash kk}), with extra wrapper packages (bm, bbm, dsfont, xspace, cancel, mathtools). The delivered numbering is the unit's own. Assets: no retained span contains a figure environment, a graphics-inclusion command, a PDF-page inclusion, a table or a TikZ picture; no image file is copied into this package, the package declares no asset dependency and no image file is redistributed with it. Licence: version arXiv:2609.24577v2 of this article is distributed under the Creative Commons Attribution 4.0 International licence, as recorded in the arXiv licence metadata for this exact version and shown on the abstract page https://arxiv.org/abs/2609.24577v2. A full rights notice is delivered at licenses/arxiv-2609.24577-CC-BY-4.0.txt. Characters: every retained excerpt is pure ASCII, so the delivered engine, XeLaTeX with its default Latin Modern text font, reports no missing character.
\begin{fact}\label{rel} 	\begin{enumerate}
\item[(i)] \(\kk[[t^2,t^3]] \;\cong\; \frac{\kk[[x,y]]}{\bigl(x^3-y^2 \bigr)}\).
\item[(ii)] 
\(\kk[[t^3,t^4,t^5]] \;\cong\; \frac{\kk[[x,y,z]]}{\bigl(y^2 - xz,\; x^3 - yz,\; z^2 - x^2 y\bigr)}\).
\item[(iii)]
\(\kk[[t^4,t^5,t^7]] \;\cong\; \frac{\kk[[x,y,z]]}{\bigl(yz - x^3,\; z^2 - xy^2,\; y^3 - x^2 z\bigr)}\).
 		\end{enumerate}

 	\end{fact}

 \begin{proposition}
 	Let
 	\(
 	R=\mathbb Z_2[[x,y,z]],
 	S=\mathbb Z_2[[t,y,z]],
 	\)
 	and define
 	\(
 	\varphi:R\to S
 	\)
 	by
 	\[
 	\varphi(x)=t^4,\qquad
 	\varphi(y)=y,\qquad
 	\varphi(z)=z.
 	\]
 	Let
 	\(
 	P=(2,\ yz-x^3,\ z^2-xy^2,\ x^2z-y^3)\subset R.
 	\)
 	Then:
 	\begin{enumerate}
 		\item $R$ and $S$ are complete regular local rings of characteristic
 		zero;
 		\item $S$ is finite free of rank $4$ over $R$;
 		\item $P$ is prime;
 	\item
 		\(
 		S/\sqrt{PS},
 		\)
 	 is regular, hence Gorenstein;
 		\item $R/P$ is not Gorenstein.
 	\end{enumerate}
 	Consequently,
 	\(
 	{
 		S/\sqrt{PS}\text{ Gorenstein}
 		\not\Longrightarrow
 		R/P\text{ Gorenstein}.
 	}
 	\)
 \end{proposition}

\begin{proof}
	
	
The ring $\mathbb Z_2$ is a complete discrete valuation ring with
 uniformizer $2$ and residue field $\mathbb F_2$.  In particular,
 $\mathbb Z_2$ is a regular local ring of dimension one and
 characteristic zero.
 Consequently
 \(
 R=\mathbb Z_2[[x,y,z]]
 \)
 and
 \(
 S=\mathbb Z_2[[t,y,z]]
 \)
 are complete regular local rings of dimension $4$ and characteristic
 zero.  Their maximal ideals are
 \(
 \mathfrak m_R=(2,x,y,z) \),
 \( \mathfrak m_S=(2,t,y,z).
 \)
The map is induced by the relation
 \(
 t^4=x.
 \)
 Thus
 \(
 S\cong R[T]/(T^4-x),
 \)
 with $T$ mapping to $t$.
 Since $T^4-x$ is monic, every element of $S$ can be written as
 \(
 a_0+a_1t+a_2t^2+a_3t^3\)
where  \(a_i\in R.
 \)
 The representation is unique, so
 \(
 S\cong R\oplus Rt\oplus Rt^2\oplus Rt^3.
 \)
  Hence
 \(
 {S\text{ is finite free of rank }4\text{ over }R.}
 \)
Consider the homomorphism
\(
\theta:\mathbb F_2[[x,y,z]]
\to
\mathbb F_2[[t]]
\)
defined by
\[
x\longmapsto t^4,\qquad
y\longmapsto t^5,\qquad
z\longmapsto t^7.
\]

The three displayed relations vanish:
\[
yz-x^3
\longmapsto
t^{12}-t^{12}=0,
\quad
z^2-xy^2
\longmapsto
t^{14}-t^{14}=0,
\quad
x^2z-y^3
\longmapsto
t^{15}-t^{15}=0.
\]

Let
\(
I:=(yz-x^3,\ z^2-xy^2,\ x^2z-y^3).
\)
Then
\(
I\subseteq\ker\theta.
\)
In view of {Fact} \ref{rel} (iii)
\(
\ker\theta=I.
\)
Since \(\mathbb F_2[[t^4,t^5,t^7]]\) is a subring of the domain
\(\mathbb F_2[[t]]\), it is a domain. 
Therefore
\(
P=(2)+I
\)
is a prime ideal of \(R\), and
\(
R/P\cong\mathbb F_2[[t^4,t^5,t^7]].
\)
Under \(x\mapsto t^4\) we obtain
\(
PS=
(2,\ yz-t^{12},\ z^2-t^4y^2,\ t^8z-y^3).
\)
Put
\(
Q=(2,y-t^5,z-t^7)\subset S.
\)
We prove
\(
\sqrt{PS}=Q.
\)
First, \(PS\subseteq Q\): modulo \(Q\) we have
\[
2=0,\qquad y=t^5,\qquad z=t^7,
\]
and hence
\[
yz-t^{12}=t^5t^7-t^{12}=0,
\quad
z^2-t^4y^2=t^{14}-t^4t^{10}=0,
\quad
t^8z-y^3=t^8t^7-t^{15}=0.
\]
Therefore
\(
PS\subseteq Q.
\)
Since \(Q\) is prime (as \(S/Q\cong\mathbb F_2[[t]]\) is a domain),
it is radical, and hence
\(
\sqrt{PS}\subseteq Q.
\)

For the reverse inclusion, work modulo \(PS\) and write
\(
\overline{S}=S/PS.
\)
Because \(2=0\) in \(\overline S\), we are in characteristic two.
The relation
\(
z^2-t^4y^2\in PS
\)
gives
\(
(z-t^2y)^2
=
z^2-2t^2yz+t^4y^2
=
z^2+t^4y^2
\in PS.
\)
Hence
\(
z-t^2y\in\sqrt{PS}.
\)
Set
$
n=z-t^2y\in\sqrt{PS},$
so 
$n^2=0 { in } S/PS.$
Then \(z=t^2y+n\). Substituting into the two remaining generators
\(yz-t^{12}\in PS\) and \(t^8z-y^3\in PS\), we obtain
\begin{align}
t^2y^2+yn &= t^{12}, \tag{1}\\
t^{10}y+t^8n &= y^3. \tag{2}
\end{align}
Multiply (1) by \(t^8\) we obtain
\(
t^{10}y^2+t^8yn=t^{20}.
\)
From (2),
\(
t^8n=y^3-t^{10}y.
\)
Substituting this into the previous equation gives
\[
t^{10}y^2+y(y^3-t^{10}y)=t^{20}
\implies
t^{10}y^2+y^4-t^{10}y^2=t^{20}
\implies
y^4=t^{20}.
\]
In characteristic two,
\(
y^4-t^{20}=(y-t^5)^4.
\)
Hence
\(
(y-t^5)^4=0
 \text{ in }S/PS,
\)
so
\(
y-t^5\in\sqrt{PS}.
\)
Since \(z-t^2y\in\sqrt{PS}\) and \(y-t^5\in\sqrt{PS}\), we get
\(
z-t^7
=
(z-t^2y)+t^2(y-t^5)
\in\sqrt{PS}.
\)
Therefore
\(
Q=(2,y-t^5,z-t^7)\subseteq\sqrt{PS}.
\)
Combining the inclusions gives
\({\sqrt{PS}=Q.}
\)
Using the above radical computation,
\[
S/\sqrt{PS}
\cong
\mathbb Z_2[[t,y,z]]/(2,y-t^5,z-t^7)
\cong
\mathbb F_2[[t]],
\]
which is regular, and every regular local ring is Gorenstein.
 We have
 \(
 R/P\cong A\cong\mathbb F_2[[t^4,t^5,t^7]].
 \)
 This is well-known that it is not Gorenstein.
 Let us show this.
 \(
 H=\langle4,5,7\rangle.
 \)
 The gaps of $H$ are
 \(
 1,2,3,6.
 \)
 Thus the genus is
 \(
 g(H)=4,
 \)
 and the Frobenius number is
 \(
 F(H)=6.
 \)
 A numerical semigroup $H$ is symmetric if and only if
 \(
 F(H)=2g(H)-1.
 \)
 Here
 \(
 F(H)=6 \)
 {but}
 \( 2g(H)-1=7.
 \)
 Hence $H$ is not symmetric.
 For a one-dimensional complete numerical semigroup ring over a field,
 \(
 \mathbb F_2[[t^H]]
 \)
 is Gorenstein if and only if $H$ is symmetric. 
 Consequently
 \(
 {R/P\text{ is not Gorenstein}.}
 \)
 \end{proof}

\end{document}
